\chapter{Pendahuluan}

\section{Tentang Dokumen Ini}

Dokumen ini adalah naskah praktikum mata kuliah \coursecode\ \coursename\ pada
Program Studi \programname, \universityname. Dokumen ini memuat sepuluh modul
laboratorium, masing-masing dirancang untuk 120 menit tatap muka ditambah tugas
mandiri di luar sesi.

Sepuluh modul dijalankan mulai minggu ke-3 sampai minggu ke-13, lalu dilanjutkan
proyek akhir pada minggu ke-14 dan ke-15. Peta lengkap modul terhadap pertemuan
kuliah dan CPMK tercantum pada bagian Panduan Praktikan.

\begin{center}
\small
\begin{tabular}{@{}p{0.26\textwidth}>{\raggedright\arraybackslash}p{0.64\textwidth}@{}}
\toprule
\textbf{Bagian} & \textbf{Isi} \\
\midrule
Panduan Asisten Praktikum & Peran asisten, alur sesi, batas bantuan, dan
  daftar periksa sebelum sesi \\
Persiapan Lingkungan & Perkakas, instalasi, uji lingkungan, aturan
  reproduksibilitas, dan penyiapan data \\
Panduan Praktikan & Ketentuan, alur sesi, format pengumpulan, peta CPMK,
  prinsip evaluasi, dan larangan \\
Modul 1--10 & Naskah praktikum, masing-masing berstruktur sama \\
\bottomrule
\end{tabular}
\end{center}

Setiap modul disusun dengan susunan yang sama: identitas, capaian,
prasyarat, konsep dasar, studi kasus, alur 120 menit, prosedur terbimbing,
latihan mandiri, tugas individu, format pengumpulan, troubleshooting, dan
ringkasan. Keseragaman ini disengaja: mahasiswa tahu di mana mencari apa tanpa
harus membaca ulang seluruh bab.

\section{Dokumen Pendamping}

Naskah ini tidak berdiri sendiri yang mana memiliki satu dokumen pendamping
dan satu repositori.

\begin{center}
\small
\begin{tabular}{@{}p{0.30\textwidth}>{\raggedright\arraybackslash}p{0.60\textwidth}@{}}
\toprule
\textbf{Dokumen} & \textbf{Isinya} \\
\midrule
\emph{\jenisdokumen\ \coursename} & Naskah ini. Konsep, prosedur, dan tugas
  untuk mahasiswa. \\
\emph{\bukupendamping} & Buku pendamping berisi rubrik penilaian tiap modul,
  bobot per kriteria, penanda jawaban yang kuat, kesalahan yang lazim muncul,
  serta lembar penilaian yang dipakai asisten. \\
\bottomrule
\end{tabular}
\end{center}

Pemisahan keduanya disengaja. Naskah praktikum dibagikan kepada mahasiswa sejak
awal semester; buku rubrik dipakai asisten saat menilai dan dibagikan kepada
mahasiswa setelah checkpoint agar mereka memahami dasar penilaiannya. Kriteria
kelulusan tiap modul tetap tercantum pada naskah ini.

\begin{catatan}
Beberapa artefak lain tidak dibagikan kepada mahasiswa sama sekali, dan itu
disengaja: daftar anomali yang ditanam pada data, kunci jawaban profiling, dan
solusi acuan tiap modul. Menemukan anomali adalah pekerjaan Modul 1 dan
Modul 10; membocorkan daftarnya menghapus seluruh nilai latihan tersebut.
\end{catatan}

\section{Repositori}

Seluruh berkas pendukung skrip SQL, pembangkit data, kerangka pipeline
Python, kerangka proyek dbt, \berkas{docker-compose.yml}, dan naskah
\LaTeX\ dokumen ini, tersedia pada repositori:

\begin{center}
\small
\fbox{\parbox{0.88\linewidth}{\centering\vspace{0.3em}
\url{\repositori}
\vspace{0.3em}}}
\end{center}

Isi repositori mengikuti susunan berikut.

\begin{center}
\small
\begin{tabular}{@{}p{0.22\textwidth}>{\raggedright\arraybackslash}p{0.68\textwidth}@{}}
\toprule
\textbf{Direktori} & \textbf{Isi} \\
\midrule
\berkas{section/} & Naskah \LaTeX\ tiap bagian dan modul \\
\berkas{sql/} & Skrip SQL per modul, termasuk \berkas{check.sql} \\
\berkas{seed/} & Skema sumber, pembangkit data, dan batch beranomali \\
\berkas{etl/} & Kerangka pipeline Python untuk Modul 7 \\
\berkas{dbt/} & Kerangka proyek dbt untuk Modul 8 dan 10 \\
\berkas{templates/} & Templat laporan, kamus data, dan dokumen desain \\
\bottomrule
\end{tabular}
\end{center}

\subsection*{Menyampaikan koreksi}

Naskah ini akan terus diperbaiki. Kekeliruan, ketidakjelasan, dan usulan
perbaikan disampaikan sebagai \emph{issue} pada repositori, disertai nomor
modul, nomor halaman, dan versi dokumen yang Anda pegang. Ketiganya
diperlukan agar koreksi dapat ditelusuri ke naskah yang tepat.

Mahasiswa yang menemukan kekeliruan pada naskah dipersilakan melaporkannya.
Laporan yang tepat dan disertai bukti diperhitungkan sebagai keaktifan.

\section{Riwayat Versi Dokumen}

Nomor versi tercantum pada halaman judul, halaman hak cipta, dan kaki setiap
halaman repositori. Sebutkan nomor ini setiap kali melaporkan kekeliruan.

\begin{center}
\small
\begin{tabular}{@{}p{0.09\textwidth}p{0.17\textwidth}>{\raggedright\arraybackslash}p{0.60\textwidth}@{}}
\toprule
\textbf{Versi} & \textbf{Tanggal} & \textbf{Perubahan} \\
\midrule
1.0 & September 2026 &
  Terbitan pertama. Sepuluh modul lengkap beserta bagian awal, skrip pemeriksa
  tiap modul, kerangka pipeline ETL dan proyek dbt, serta pembangkit data
  NusaMart tiga ukuran dan tiga sumber. \\
\bottomrule
\end{tabular}
\end{center}

\begin{catatan}
Setiap perubahan naskah yang berarti menaikkan nomor versi dan menambah satu
baris pada tabel di atas. Angka pada berkas \berkas{config/metadata.tex} dan
tabel ini harus selalu sejalan --- naskah yang nomor versinya tidak dapat
dipercaya membuat pelaporan kekeliruan menjadi mustahil ditelusuri.
\end{catatan}

\section{Untuk Siapa Dokumen Ini}

\begin{center}
\small
\begin{tabular}{@{}p{0.20\textwidth}>{\raggedright\arraybackslash}p{0.70\textwidth}@{}}
\toprule
\textbf{Pembaca} & \textbf{Mulai dari} \\
\midrule
Mahasiswa & Panduan Praktikan, lalu modul sesuai pekan berjalan. Kerjakan
  bagian Pre-lab \emph{sebelum} sesi 120 menit tatap muka tidak cukup bila
  persiapan dikerjakan di kelas. \\
Asisten & Panduan Asisten Praktikum, lalu \berkas{seed/README.md} pada
  repositori untuk penyiapan data. \\
Pengajar & Pendahuluan ini, lalu peta modul dan CPMK pada Panduan Praktikan. \\
\bottomrule
\end{tabular}
\end{center}

Bagi mahasiswa, satu saran praktis: bacalah bagian Ringkasan pada akhir setiap
modul \emph{sebelum} mengerjakan modul tersebut. Ringkasan menyebutkan tiga
gagasan yang akan dibawa ke modul berikutnya, dan mengetahui tujuannya lebih
dahulu membuat prosedur terbimbing lebih mudah diikuti.

\clearpage
